\documentclass[../main]{subfiles}
\begin{document}
\ifSubfilesClassLoaded{\setcounter{page}{54}}{}
\section{Коммунизм и общественное сознание}
\label{sec:10_kommunizm_soznanie}

С точки зрения марксизма, совершая материальное производство общество производит своё собственное сознание. Человечество знает различные исторически сложившиеся формы общественного сознания. Из мифологического сознания, где идеальное ещё не отделено, или на разных стадиях не вполне отделено от непосредственной материальной деятельности, с развитием разделения труда и отчуждения человека от своей собственной деятельности выделяются философия и религия, затем искусство и наука. Наряду с ними существуют мораль и право, выделившиеся тогда, когда общество стало нуждаться во внешнем регулировании взаимоотношений между людьми с противоположно направленными интересами. Общественное сознание — это не сумма сознаний отдельных индивидов, а сознание общества как целого, отделившееся от непосредственной человеческой деятельности, её идеальное измерение. И наоборот, индивидуальное сознание — это то же общественное сознание, приобретенное человеком в зависимости от его места и роли в общественной жизни.

Общественное сознание и его формы обладают по отношению к членам общества объективным существованием. Чтобы жить в человеческом обществе, с требованиями и установками морали, права, религии, науки и даже искусства и философии люди вынуждены считаться точно так же, как и с окружающей природой и уровнем развития технологий. В общественном сознании отражено и выражено не нечто произвольное, а способ взаимодействия человека с природой и человека с человеком. Общественное сознание и все его формы - не непосредственный продукт природы, а продукт её освоения человеческим обществом, отражённый способ производства.

Факт отделения общественного сознания от непосредственной деятельности и приобретение его формами относительно-самостоятельного существования, а также его принудительный обязательный характер, по отношению к членам общества, был зафиксирован задолго до Маркса. Он лежит в основе не только религиозного миропонимания, но и всей идеалистической философии.

Становление идеи коммунизма требовало критики как самого этого факта, так и того, как он интерпретируется в теоретическом мышлении. Критика общественного сознания не просто своей эпохи, а общества, где господствует эксплуатация человека человеком вообще, сыграла важнейшую роль в становлении идейного содержания коммунизма.

Марксизм, как известно, прошёл школу критики философии Гегеля и Фейербаховской критики христианской религии - формы общественного сознания, которая была господствующей в Европе на протяжении Средних Веков и продолжает претендовать на господство и в буржуазном обществе. До какой степени серьезно Маркс относился к фейербаховской критике религии видно из того, что он говорит ни больше ни меньше как о подвиге Фейербаха. В чем же по отношению к нашему вопросу, согласно Марксу, состоит этот подвиг? «1) в доказательстве того, что философия есть не что иное, как выраженная в мыслях и логически систематизированная религия, не что иное, как другая форма, другой способ существования отчуждения человеческой сущности, и что, следовательно, она также подлежит осуждению (это - очень важное замечание, относящиеся к философии как к форме общественного сознания); 2) в основании истинного материализма и реальной науки, поскольку общественное отношение "человека к человеку" Фейербах также делает основным принципом теории».

Это достоинства концепции Фейербаха, но она не до конца последовательна, и Маркс специально отметил это: «Фейербах исходит из факта религиозного самоотчуждения, из удвоения мира на религиозный, воображаемый мир и действительный мир. И он занят тем, что сводит религиозный мир к его земной основе. Он не замечает, что после выполнения этой работы главное-то остается еще не сделанным. А именно, то обстоятельство, что земная основа отделяет себя от самой себя и переносит себя в облака как некое самостоятельное царство, может быть объяснено только саморазорванностью и самопротиворечивостью этой земной основы. Следовательно, последняя, во-первых, сама должна быть понята в своем противоречии, а затем практически революционизирована путем устранения этого противоречия. Следовательно, после того как, например, в земной семье найдена разгадка тайны святого семейства, земная семья должна сама быть подвергнута     теоретической    критике     и   практически      революционно преобразована».

Земная жизнь человека - материальное производство - таким образом, становится для Маркса тем предметом, который стоит изучать, чтобы понять сущность земных корней сознания. «Религия сама по себе лишена содержания, её истоки находятся не на небе, а на земле, и с уничтожением той извращенной реальности, теоретическим выражением которой она является, она гибнет сама собой». (Маркс — Арнольду Руге в Дрезден (30 ноября 1842 г.)

То же самое относится и к другим формам общественного сознания: уже упомянутая философия, мораль, право, и даже наука и искусство, должны быть поняты как стороны человеческой практики из неё самой, из её характера и содержания, а не из самих себя. В марксизме чётко разделяется базис общественных отношений - единство производительных сил и производственных отношений - и производная от них надстройка, к которой относится общественное сознание.

Например, о праве Маркс в подготовительных рукописях к капиталу пишет следующее: «Если эти тривиальности (представление о праве современных Марксу экономистов - ред.) свести к их действительному содержанию, то они скажут больше, чем знают их проповедники. А именно, что каждая форма производства порождает свойственные ей правовые отношения, формы правления и т. д. Грубость и отсутствие понимания в том и заключается, что органически между собой связанные явления ставятся в случайные взаимоотношения и в чисто рассудочную связь. Буржуазным экономистам мерещится только, что при современной полиции можно лучше производить, чем, например, при кулачном праве. Они забывают только, что и кулачное право есть право и что право сильного в другой форме продолжает существовать также и в их «правовом государстве».

«Земные корни» общественного сознания должны быть не только поняты, но и изменены - изменен характер общественного производства. И то, что формы общественного сознания произрастают из способа общественного производства, свидетельствует о том, что с ними дела обстоят вовсе не так просто, как может показаться на первый взгляд. Например, полностью снять право, преодолеть необходимость его как внешнего регулятора можно только в определённых общественных условиях - при коммунизме, развитом на своём собственном основании. То же самое относится и к преодолению религии. Её нельзя просто отменить, объявив ложной. Её нужно снять, сохранив все достижения человеческой культуры, которые развивались в религиозной форме, уничтожив путем коренного изменения характера ежедневно совершающейся человеческой деятельности.

Конечно, в маленькой брошюрке об общественном сознании при коммунизме мы можем говорить только в общих чертах. Но даже в этом случае следует отметить одно важнейшее обстоятельство, относящееся к предшествующему коммунизму обществу, разделённому на классы. Точно так же, как люди занимают разное положение в процессе производства, соответствующее их классовой принадлежности, так и общественное сознание носит классовый характер. Это не значит, что не существует единого поля сознания и объективных, независимых только и исключительно от классовой позиции её носителя, истины, добра и красоты. Это значит, что само это поле сознания, является полем классовой борьбы. Это значит, что господствующий экономический класс господствует также и в сознании. Что сознание господствующего класса является господствующим в классовом обществе. Поэтому Ленин вслед за Марксом и Энгельсом говорит о важности борьбы на поле сознания – о теоретической борьбе, как об особом виде классовой борьбы.

В брошюре «Что делать?» Ленин посвятил целый пункт значению теоретической борьбы пролетариата и точке зрения Энгельса на этот вопрос. Очень актуальными на сегодняшний день нам представляется следующее замечание Ленина: «Кто сколько-нибудь знаком с фактическим состоянием нашего движения, тот не может не видеть, что широкое распространение марксизма сопровождалось некоторым принижением теоретического уровня. К движению, ради его практического значения и практических успехов, примыкало немало людей, очень мало и даже вовсе не подготовленных теоретически. Можно судить поэтому, какое отсутствие такта проявляет «Раб. Дело», когда выдвигает с победоносным видом изречение Маркса: «каждый шаг действительного движения важнее дюжины программ». Повторять эти слова в эпоху теоретического разброда, это все равно что кричать «таскать вам не перетаскать!» при виде похоронной процессии. Да и взяты эти слова Маркса из его письма о Готской программе, в котором он резко порицает допущенный эклектизм в формулировке принципов: если уже надо было соединяться — писал Маркс вожакам партии — то заключайте договоры, ради удовлетворения практических целей движения, но не допускайте торгашества принципами, не делайте теоретических «уступок». Вот какова была мысль Маркса, а у нас находятся люди, которые, во имя его, стараются ослабить значение теории! Без революционной теории не может быть и революционного движения. Нельзя достаточно настаивать на этой мысли в такое время, когда с модной проповедью оппортунизма обнимается увлечение самыми узкими формами практической деятельности.»

Обращаясь к наследию Энгельса, Ленин приводит большую цитату, которую мы считаем нужным также привести и здесь: «Немецкие рабочие имеют два существенных преимущества перед рабочими остальной Европы. Первое — то, что они принадлежат к наиболее теоретическому народу Европы и что они сохранили в себе тот теоретический смысл, который почти совершенно утрачен так называемыми «образованными» классами в Германии. Без предшествующей ему немецкой философии, в особенности философии Гегеля, никогда не создался бы немецкий научный социализм, — единственный научный социализм, который когда-либо существовал. Без теоретического смысла у рабочих этот научный социализм никогда не вошел бы до такой степени в их плоть и кровь, как это мы видим теперь. А как необъятно велико это преимущество, это показывает, с одной стороны, то равнодушие ко всякой теории, которое является одной из главных причин того, почему английское рабочее движение так медленно двигается вперед, несмотря на великолепную организацию отдельных ремёсел, — а с другой стороны, это показывает та смута и те шатания, которые посеял прудонизм, в его первоначальной форме у французов и бельгийцев, в его карикатурной, Бакуниным приданной, форме — у испанцев и итальянцев.

Второе преимущество состоит в том, что немцы приняли участие в рабочем движении почти что позже всех. Как немецкий теоретический социализм никогда не забудет, что он стоит на плечах Сен-Симона, Фурье и Оуэна — трех мыслителей, которые, несмотря на всю фантастичность и весь утопизм их учений, принадлежат к величайшим умам всех времен и которые гениально предвосхитили бесчисленное множество таких истин, правильность которых мы доказываем теперь научно, — так немецкое практическое рабочее движение не должно никогда забывать, что оно развилось на плечах английского и французского движения, что оно имело возможность просто обратить себе на пользу их дорого купленный опыт, избежать теперь их ошибок, которых тогда в большинстве случаев нельзя было избежать. Где были бы мы теперь без образца английских тред-юнионов и французской политической борьбы рабочих, без того колоссального толчка, который дала в особенности Парижская Коммуна?

Надо отдать справедливость немецким рабочим, что они с редким уменьем воспользовались выгодами своего положения. Впервые с тех пор, как существует рабочее движение, борьба ведется планомерно во всех трех ее направлениях, согласованных и связанных между собой: в теоретическом, политическом и практически-экономическом (сопротивление капиталистам).»

Это место показывает отношение Энгельса к теоретической борьбе. Но классовая борьба в области общественного сознания не исчерпывается одной теоретической борьбой, точно так же, как и общественное сознание не ограничивается понятийным мышлением. Борьба ведётся и в области эстетической деятельности, и в области морали, и по отношению к праву и религии. Этого тоже нельзя недооценивать. Однако, теоретическая работа по приобретению и развитию пролетариатом своего собственного классового сознания является ключом ко всем остальным формам сознания, поскольку позволяет их понять.

Важно также отметить, что в классовом обществе в общественном сознании, вне зависимости от того понимают это люди или нет, общечеловеческое реализуется через классовое. Сознание революционного класса, выражающего на данном историческом этапе общечеловеческий интерес, и является движущей силой истории, и реализует это общечеловеческое. В этом смысле не только мораль, религия и философия, право, но и наука и искусство являются также и идеологией, поскольку и эстетическое освоением мира в искусстве, и познание в науке, обусловлено, а значит и ограничено, (свобода, которая присутствует в нем, ограничена) историческим движением общества и расстановкой сил в классовой борьбе.

В связи с этим хотелось бы обратиться к решению проблемы о субъекте мышления Марксом и Лениным. Нужно напомнить о том, что Маркс делил политическую экономию на классическую и вульгарную, избрав для этого совершенно определенный критерий. Он указывает на то, что классическая политэкономия была этапом развития понимания своего предмета, в отличии от вульгарной. Он четко определяет партийность, то есть классовость этой науки. И классическая, и вульгарная политическая экономия - наука буржуазная, стоящая на службе у буржуазии и выражающая в теории направление ее развития. Как только это развитие перестает совпадать с развитием общества, перестает воплощать это развитие, другими словами, буржуазия перестает быть революционным классом, политическая экономия становится вульгарной. Это означает, что она уже не является развивающейся мыслью, а наоборот становится преградой на пути развития мысли.

Разумеется, дело тут не в личных качествах отдельных политэкономов. Маркс убедительно показывает, что субъектом мышления являются не отдельные теоретики и не общество в целом, а именно класс. Более того, он является субъектом мышления только до тех пор, до каких он является субъектом исторического действия.

Ту же мысль мы находим и у Ленина, когда он говорит о партиях в философии, которых может быть только две: мыслящая партия, и партия, препятствующая мышлению. И именно потому, что одна партия - это партия класса, который является субъектом исторического действия, а значит, и субъектом мышления, а другая не является таковой. Это касается не только философии, но ВСЕХ без исключения наук, в том числе и естественных. Хотя там это не так очевидно, но, например, получить финансирование не приносящих даже в перспективе плодов, исследования становится настолько трудно, что это вполне можно рассматривать как эпизоды завуалированной классовой борьбы. Это остается актуальным до тех пор, пока развитие общества осуществляется через борьбу классов. Получается, что объективная истина всегда тенденциозна, имеет классовый характер. То же относится к объективности добра и красоты.

Ленин в сентябре (августе по старому стилю) 1917 года высказал хорошо известную и, на наш взгляд, очень важную мысль в духе написанного выше, показывающую, как в революционные эпохи пролетариат относится к своему классовому сознанию, учитывая, что цель пролетариата, в отличии от других классов заключается вовсе не в том, чтобы завоевать и упрочить свое классовое господство, а в том, чтобы завоевав это господство уничтожить всякое классовое господство и всякое деление общества на классы вообще: «Будем непреклонны,» - говорит Ленин, - «в разборе малейших сомнений судом сознательных рабочих, судом своей партии, ей мы верим, в ней мы видим ум, честь и совесть нашей эпохи, в международном союзе революционных интернационалистов видим мы единственный залог освободительного движения рабочего класса». «Господство класса невозможно без господства его в области сознания», - подчеркивает Ленин. Чтобы победить, революционный класс в области сознания должен мерить себя не чужой мерой, а своей, определяемой его исторической миссией. Вырвав у буржуазии господство в этой области, вооружившись передовыми достижениями своего собственного классового сознания, пролетариат становится действительной материальной силой исторических преобразований: «Оружие критики не может, конечно, заменить критики оружием, материальная сила должна быть опрокинута материальной же силой, но теория становится материальной силой, как только она овладевает массами», - писал Маркс в «К критике гегелевской философии права».

Однако, до сих пор мы говорили об общественном сознании не при коммунизме, а о положении дел от которого нужно отталкиваться при переходе к нему. На чём же по отношению к общественному сознанию настаивает коммунизм как наука?

Коммунизм настаивает на преодолении общественного разделения труда, отчужденности человека от других людей, от процесса и результата своей деятельности. А это влечёт за собой и соответствующие изменения в области общественного сознания. Отчуждённость форм общественного сознания как друг от друга, так и непосредственно от человеческой деятельности, тоже будет преодолена.

Идеальные деятельные способности человека, которые «замкнуты» в рамках форм общественного сознания, как отделённых от материального производства сферах эстетической, нравственной и интеллектуальной деятельности, станут достоянием всех и каждого.

Вопрос об общественном сознании при коммунизме - это вопрос практики. Это вопрос планирования и реализации запланированного в материальной деятельности так, чтобы происходили соответствующие изменения и в духовной сфере. Изменяя свою деятельность, люди изменяют и своё сознание как идеальную сторону этой деятельности. Но само изменение практики носит при коммунизме сознательный характер. Потому здесь уже не просто общественное бытие определяет общественное сознание, а общественное сознание, как сознание согласованно-действующей свободной ассоциации людей, определяет бытие, определяющее это сознание.

Обобществление всей духовной культуры, причем по логике этой культуры, становится острейшей необходимостью для перехода к коммунизму. Поэтому и в теории коммунизма проблемы общественного сознания имеют смысл только при условии, если понимать их как вопросы становления коммунизма.

В коммунистической деятельности, в отличии от отчуждённого труда, идеальное как соответствующее идеалу и идеальное как относящееся к «миру идей», в конечном счёте (именно в конечном счёте), эмпирически совпадают. Идеальное, таким образом, выступает и как отражение (представление, выражение), материального в материальном - деятельное отражение объективной реальности, и как идеальная цель деятельности, с которой она должна сообразовываться, и как эстетический идеал. В этом заключается важнейший момент преодоления разделения труда, его отчуждённости во всех смыслах слова, и утверждается самоценность деятельности вне зависимости от того, что именно делается.

Здесь следует отметить, что неверно было бы утверждать, что сознательный и материальный моменты человеческой деятельности равноправны в труде как таковом, вне зависимости от его исторической специфики, и уж тем более, на уровне деятельности общества как целого. Конечно, идеальный момент присутствует в любой деятельности как момент целеполагания, если только целеполагание не лежит за пределами этой деятельности (для работника, который только продаёт рабочую силу, его собственная цель вовсе не совпадает с целью производства), но речь идёт о другом. В любом случае, способ производства материальной жизни определяет способ производства идеального. И с этим приходится считаться.

Это важно в практическом плане, потому, что, хотя идея и становится материальной силой, когда она овладевает массами, сама она есть результат осознания материальных сил в процессе их действия (результат и сторона взаимодействия материального с материальным). Другими словами, чтобы изменить способ производства идей, нужно изменить способ производства материальной жизни. Только так, а не наоборот. Наоборот ничего не выйдет.

Например, сколько бы мы сейчас не пропагандировали (то есть, действовали идеально) диалектику, она не станет от этого обыденным мышлением, а производство диалектического мышления пока так и будет кустарным, если можно так выразиться, делом. Эта кустарщина, конечно, очень важна, и от неё во многом зависит будущее. Но мы должны понимать, что диалектически мыслящие люди не будут массовым продуктом при таком, как сейчас, способе производства материальной жизни и соответствующем ему способе производства сознания.

Для проблемы коммунизма, нужно не просто зафиксировать направление перехода материальной жизни человека в сознание (что первично, а что вторично), но и выдвинуть идеи по практическому совершению этого перехода. И тут очень важны частные идеи на этот счёт. Важны идеи по формированию тех или иных эмпирических форм коммунизма, так как никогда всё материальное в одночасье не переходит в сознание, а всегда делает это в каких-то особенных формах. Важны, например, такие идеи, как у Макаренко: любовь - это обычное дело, которое нужно правильно организовать.

Какую бы огромную роль ни играл идеальный момент деятельности, речь всегда идёт о жизни и деятельности материального человека в материальном мире. Идеальное всегда является атрибутом материальной деятельности общественного человека и соответствует её движению. «Царство свободы», тотальный коммунизм, характеризуется в том числе и сознательностью производства общественного сознания в общественном же масштабе.

Отдельные прорывы в сфере общественного сознания, выходящие за рамки капитализма, происходят уже сейчас. Речь о них (в том числе) шла во второй главе этой работы. Это производство идеального при создании современных технологий. Но дело не только в этом, не всякое идеальное относится к общественному сознанию. Это уже не просто идеальная деятельность в своей отчуждённости, которая носит надстроечный характер, и не замкнутое в рамках профессионального кретинизма идеальное. Это идеальный момент производства человеческой жизни вообще. И то, как именно эта жизнь производится, определяет то, (как и в случае с базами данных и со свободным ПО, имеет место один и тот же процесс) как создается новое непосредственное, неотчуждённое общественное сознание, пусть пока на локальных участках. Это такое сознание, которое производится не только и не столько самим фактом производства программы, сколько самим характером её производства и функционирования.

Общественное сознание может определять общественное сознание, только определяя общественное бытие. Причём последнее нельзя понимать, как нечто отдельное от целостного человека, как что-то идеальное вообще, а только как момент деятельности материального человека в материальном же мире. Дело заключается именно в сознательном производстве целостного человека - универсально действующего, чувствующего, думающего. Такая деятельность вне зависимости от того, что именно делается, будет выступать как осуществление коммунизма в непосредственно коллективной форме и будет деятельностью по производству свободы, целеполагания, чувственности, нравственности и мышления.

По большому счёту, это предполагает, что любая деятельность при коммунизме должна в том числе быть производством идей и идеала будущей жизни, производством человека.
\end{document}
